\documentclass[brazil, english, a4paper, 12pt]{article}

\usepackage{abakos} 
\newcommand{\monog}{}
\title{Os Três Tipos Lógicos De Pesquisa}
\author{Gustavo Martins de Avila}
\date{14 de Abril de 2026}

\begin{document}

\maketitle

\noindent \textbf{Resumo:} Esse trabalho sintetiza os conceitos fundamentais apresentados por Gilson Volpato sobre a logica da investigação científica. Discutem-se as distinções entre Descrição, Associação e Associação com interferência, ressaltando o papel da hipótese em cada modelo.

\vspace{0.5cm}

\section{OBJETIVO}
Este artigo tem como objetivo explicar os conceitos dos três tipos logicos de pesquisa: Descrição, Associação e Associação com interferência, visando compreender as motivações e o desenvolvimento de um artigo científico.

\section{DISCUSSÃO}
A lógica da pesquisa científica é determinada pela pergunta inicial. Enquanto pesquisas descritivas focam em caracterizar uma variável e dispensam hipóteses, as pesquisas de associação buscam relações de covariância. Na associação com interferência, a hipótese é o motor, propondo um mecanismo de causa e efeito a ser testado.

\section{CONCLUSÃO}
A clareza sobre o tipo lógico da pesquisa é essencial para a eficiência do método. A escolha correta evita erros teóricos e direciona a coleta de dados e a argumentação necessária para o avanço científico.

\vspace{1cm}

\renewcommand\refname{REFERÊNCIAS}
\begin{thebibliography}{9}

\bibitem{volpato} 
VOLPATO, Gilson. \textbf{Os três tipos lógicos de pesquisa}. Publicado pelo canal Gilson Volpato. 15 mar. 2012. Disponível em: <http://www.youtube.com/watch?v=XoTQo7fUfos>. Acesso em: 14 abr. 2026.

\bibitem{puc} 
PUC MINAS. \textbf{ABNT: Elaborar e formatar artigo científico}. Belo Horizonte: Biblioteca PUC Minas, 2024. Disponível em: <https://www.pucminas.br/biblioteca>. Acesso em: 14 abr. 2026.

\end{thebibliography}

\end{document}